%%%%%
%
\section{R\'efractivit\'e 1\`ere version}

\itemstyle{Op\'erateur~:} {\sf IndRefra}
\index{IndRefra}

%%%%%OBJET%%%%
\itemstyle{Objet~:}
Calcul de la r\'efractivit\'e de l'air \`a 
partir de la mesure du r\'efractom\`etre.

%%%%%ARGUMENTS%%%%
\bigskip
%%%%%ENTREES%%%%
\begin{tabarg}
\multicolumn{3}{l}{\titretab {Entr\'ees}}\\
\hline
$V$ & Vecteur, {\tt float} & Tension de sortie de l'appareil ($V$)\\
$V_0$ & Coeff., {\tt float} & Coefficient li\'e \`a l'appareil\\
$K$ & Coeff., {\tt float} & Coefficient li\'e \`a l'appareil\\
$N_0$ & Coeff., {\tt float} & Coefficient li\'e \`a l'appareil\\
$F_0$ & Coeff., {\tt float} & Coefficient li\'e \`a l'appareil\\
\end{tabarg}

\medskip
%%%%%SORTIES%%%%
\begin{tabarg}
\multicolumn{3}{l}{\titretab {Sorties}}\\
\hline
$N$ & Vecteur, {\tt float} & R\'efractivit\'e\\
\end{tabarg}


%%%%%FORMULE%%%%
\itemstyle{Formule}

%%IndRefra
\[
N = N_0 + \frac{10^6 - N_0}{9450 - F_0}
    .\frac{V - V_0}{K}
\]


%%%% COMMENTAIRE %%%%
\itemstyle {Commentaire~:}

Formule employ\'ee jusqu'en 1994 inclus, avec les valeurs suivantes :\\
$V_O = 0.341$, $K=1.235$, $N_0= 244.5$, $F_0=4.385$.


%%%%%AUTEUR%%%%
\itemstyle{Auteur~:}
CNRM/GMEI/TRAMM

%%%%%REFERENCE%%%%
\itemstyle{R\'ef\'erences~:}

Obtenu aupr\`es de J-Y. Delahaye (CETP).
